\section{The argument as a whole}\label{s:whole}

The commentary on the lemmata has examined the text sentence by sentence. This section considers it as a whole: the structure of the argument, the readings of the whole that the text admits, the objections that are not objections to any one lemma, and what the text, with the Commentary, does not establish.

\subsection{Structure}

The argument of EXI-01-a, as the Commentary has reconstructed it, has the following structure.

\begin{enumerate}[label=(\roman*)]
\item Definition~1 says what meaning is: the extent to which a life consists in active engagement with projects of objective worth. It is a real definition, defended by Wolf, read conatively in its subjective component (\Lm{2.5}), broadly in its notion of a project (Proposition~\ref{p:broad}), and with a measure left unspecified but constrained (Propositions~\ref{p:noneg} and~\ref{p:neutral}).
\item Assumption~2 says what objective worth is, in a restricted sense (Proposition~\ref{p:restricted}), and that there is some. It plays no role in the derivation (Proposition~\ref{p:redundant}) and is independent of the other premises (Proposition~\ref{p:indep}); its role is to interpret the theorem and to fix the text's position on the sceptical challenge. It is defended in Section~\ref{s:defence} and entailed by the theory of value of ETH-01 (Proposition~\ref{p:eth01}).
\item Assumption~3 says that every agent capable of intentional action can engage in some project of objective worth. It carries the weight of the availability claim, it is derived in EXI-01-c from a premise about the reach of worthwhile objects, and it is supported by EXI-17-a's demonstration that meaning does not require success.
\item The bridge principle says that what an agent can engage in, some life available to her includes. It is unstated in the text, explicit in the Lean module, and analytic on the reading of ability adopted here (Proposition~\ref{p:bridge}).
\item Theorem~4 follows from (i), (iii) and (iv): a life of positive meaning is available to every agent, at every time at which she is one (\Lm{5.3}).
\item Remark~5 says what the theorem answers and what it does not. The sceptical challenge is not answered; its vindication premise is denied by Assumption~2 (Proposition~\ref{p:vind}). The absurd in Nagel's sense survives and is compatible with meaning (Proposition~\ref{p:compat}).
\end{enumerate}

\subsection{Readings of the whole}

Three readings of the text as a whole should be distinguished.

\paragraph{The deflationary reading.} On this reading, EXI-01-a is a trivial text. Its theorem follows in three lines from assumptions that say almost the same thing, and the module's header concedes that the proof is immediate. The text adds nothing to Wolf except the claim of availability, and that claim is purchased by an assumption that simply asserts it.

The reading is mistaken, though it starts from a truth. It is true that the derivation is short. But the shortness of a derivation is not the triviality of a text, and the deflationary reading overlooks three things the text does. It separates the availability claim from the definition of meaning, which Wolf's lectures do not; it isolates the assumption on which availability depends, which makes the claim testable and shows exactly what a critic must deny; and it locates its own position on the sceptical challenge with a precision that the literature on Nagel rarely achieves. The comments on the lemmata have shown that each of these is substantive.

\paragraph{The programmatic reading.} On this reading, EXI-01-a is the statement of a position whose defence is distributed: the definition is defended by Wolf, the existence of objective worth by Metz and by ETH-01, the availability claim by EXI-01-c and EXI-17-a, its consequences by EXI-01-d, its relation to the absurd by EXI-04 and by this Commentary, its relation to death by EXI-02, EXI-03 and EXI-19, its relation to theism by EXI-11. The page states the position in its barest form, and the family and the discipline supply the rest.

We think this is the correct reading. It explains the text's form: an axiomatic statement is the natural form for a position whose parts are defended elsewhere, since it shows how the parts fit together. It explains the citations, which are not of authorities for the claims but of the places where the claims are defended. And it explains Remark~5, which reads as a guide to where the defence of the text is to be found.

\paragraph{The dialectical reading.} On this reading, the purpose of EXI-01-a is to fix the terms of debate: to say what a theory of meaning must contain (a definition, an assumption about value, an assumption about reach) and to show where each rival theory disagrees. The subjectivist denies Assumption~2's gloss; the supernaturalist denies that Assumption~2 can be true without God; the sceptic of the absurd denies Assumption~2 itself, or demands its vindication; the pessimist denies Assumption~3 for some agents. The reading is compatible with the programmatic one, and the Commentary has used it in identifying, at each lemma, the rival positions that deny it.

\subsection{The family read with the lemmata}\label{s:family}

The two other manuscripts of the family were written after EXI-01-a and with it in view. Read with the lemmata, they show which parts of the text the family regarded as needing support and how it supplied it.

\paragraph{EXI-01-c.} The alternative proof keeps Definition~1 and Assumption~2 and replaces Assumption~3 (\Lm{4}). Its Definition~2.1 introduces objects of engagement, which include persons, living creatures, practices, inquiries and parts of the world as objects of understanding, and calls them worthwhile if they have or serve attitude-independent value. Its Assumption~2.2 (Reach) says that every agent capable of intentional action can direct caring activity at such an object. Its Lemma~2.3 says that sustained caring activity directed at a worthwhile object is a project of objective worth in the sense of Definition~1. Its Theorem~2.4 is Theorem~4.

Three lemmata of EXI-01-a are affected. At \Lm{2.6}, EXI-01-c's Lemma~2.3 is the argument for the broad reading of \enquote{projects} (Proposition~\ref{p:broad}). At \Lm{4.3}, the quantifier over projects is replaced by a quantifier over objects, and Assumption~3 becomes a corollary. At \Lm{4.4}, \enquote{able to engage} is analyzed as \enquote{able to direct caring activity at}, which is the conative reading of engagement defended at \Lm{2.5}. The alternative proof does not affect the bridge principle (Proposition~\ref{p:bridge}), which it also uses without stating, or the role of Assumption~2, which it also does not use in its derivation. The Lean module could be adapted to it by replacing the axiom \texttt{assumption\_3} with two axioms, one for Reach and one for the lemma, and deriving \texttt{assumption\_3} as a theorem; the axiom list of \texttt{meaning\_available} would then contain Reach, the lemma and the bridge.

\paragraph{EXI-01-d.} The consequence manuscript draws four propositions from the availability claim. Its Proposition~2.1 (independence from outcome) is EXI-17-a's Corollary~3.6, and bears on \Lm{5.3}. Its Proposition~2.2 (independence from talent and scale) bears on \Lm{2.7}: it uses the degree reading of objective worth, on which scale affects degree but not presence. Its Proposition~2.3 (no condemnation by circumstance) is Theorem~4 restated as a denial, with the strengthening of EXI-01-c. Its Proposition~2.4 (degree, not possibility) bears on \Lm{2.3} and \Lm{10}: circumstances affect the degree of meaning attainable but not its availability. Its section on what availability does not imply corresponds to Remark~5's second sentence (\Lm{10}). Its reply on the absurd repeats Remark~5's third sentence and cites EXI-04.

\paragraph{The division of labor.} Read together, the three manuscripts divide the family conclusion as follows. Its first half, that life is meaningful to the extent that one engages with projects of objective worth, is Definition~1, with its objective component interpreted by Assumption~2 and its project component interpreted by EXI-01-c's Lemma~2.3. Its second half, that meaning so understood is available to everyone, is Theorem~4, derived in EXI-01-a from Assumption~3 and in EXI-01-c from the weaker-seeming Reach, and its consequences and limits are drawn in EXI-01-d. The present work supplies what Remark~5 says is missing: the defence of Assumption~2 and the discussion of the absurd.

\subsection{Objections to the whole}

\begin{objection}[Circularity]\label{o:circ}
The availability of meaning is secured by stipulation. Definition~1 is tailored so that the conditions of meaning are within everyone's reach, and Assumption~3 then says that they are. A theory of meaning should be tested against our judgments of which lives are meaningful, not designed to deliver a desired distribution.
\end{objection}

\begin{reply}
Definition~1 was not designed to deliver availability. It is Wolf's definition, defended by her on the basis of judgments about cases, and Wolf does not argue for availability. The availability claim is a consequence, derived in the family by a separate argument (EXI-01-c) from a separate premise. That the definition, so supported, has this consequence is a discovery about it, and a welcome one; it is not a stipulation. The test of the definition against cases is the one the comments on \Lm{2.5} and \Lm{2.7} have applied, and it is passed independently of availability.
\end{reply}

\begin{objection}[Meaning and morality]\label{o:gauguin}
Williams's painter abandons his family to pursue his art \citep{Williams1981b}. On Definition~1, his life gains meaning from his engagement with painting, which has objective worth, while the wrong he does his family is irrelevant to meaning. A theory on which a life can be made more meaningful by a moral wrong is defective.
\end{objection}

\begin{reply}
The objection correctly reports what Definition~1 implies, and the implication is correct. Wolf argued that meaning is a value distinct from morality, and the painter's case is a case in which the two come apart: his life gains meaning from his engagement with painting and is morally worse for the abandonment. To say that his life is more meaningful is not to say that it is better all things considered, or that the abandonment was justified, or that anyone should admire it. A theory that refused to recognize the meaning of his engagement because of the wrong would be confusing meaning with morality. The case also illustrates EXI-17-a's thesis. Williams argued that whether the painter's choice was justified depends on whether he succeeds; on Definition~1, the meaning of his life does not depend on it. What his success affects is the justification of his choice and the value his work adds to the world, not the extent to which his life consisted in engagement with a worthwhile project.
\end{reply}

\begin{objection}[Shape]\label{o:shape}
The meaning of a life depends on its shape: on whether it progresses, whether its later parts redeem its earlier failures, whether it ends well \citep{Velleman1991,Kauppinen2012}. Definition~1 measures meaning by extent, which is indifferent to order, and so cannot be the right account.
\end{objection}

\begin{reply}
Definition~1 does not say that extent is indifferent to order; the comment on \Lm{2.3} has argued that the measure is left open, and Proposition~\ref{p:neutral} that the availability claim is independent of it. EXI-19 holds that the value of a life as a whole depends in part on its shape (Principle~P1 of EXI-19-a), and EXI-19-b that finitude is a condition of the meaning of a life as a whole, which depends on its shape (its Proposition~1.2). The view that shape matters to meaning is therefore part of the discipline, and it is compatible with Definition~1 on any measure on which the extent to which a life consists in engagement is assessed over the whole and not summed over moments. What would be incompatible with Definition~1 is the view that shape could confer meaning on a life that contains no engagement with anything of worth, and that view has no defenders: a life that progresses steadily toward the counting of more blades of grass has a shape but no meaning.
\end{reply}

\begin{objection}[Transcendence]\label{o:transc}
\citet{Nozick1981} argued that a life is meaningful only through connection with something beyond itself, and that the question of meaning recurs for whatever it connects with until it reaches something unlimited. Definition~1 stops the regress arbitrarily at projects of objective worth.
\end{objection}

\begin{reply}
The regress arises only if the worth of what one connects with must itself be explained by connection with something further. On Assumption~2, the worth of a project is not conferred by its connection with anything; it belongs to the project, or to the object it serves, in virtue of what it is. The regress stops where worth is found, and it is found in the care of persons, in understanding and in the other components of the good. EXI-17-a accepts that meaning requires connection with value beyond oneself (its reply to Objection~2), and Definition~1 secures it through the objective component. What Nozick's regress adds is the demand that the endpoint be unlimited, and that is a demand for the kind of meaning that EXI-11-a says a theistic framework would add (its Proposition~4.1), not a condition on meaning as such.
\end{reply}

\begin{objection}[Theism]\label{o:theism}
If REL-01 is right that a maximally great being exists, the meaning of a life consists in its relation to that being, and Definition~1 describes at best a secondary meaning.
\end{objection}

\begin{reply}
EXI-11 holds that meaning does not require God, though a theistic framework would add to it. The addition is accommodated by Definition~1's measure: engagement with a being of maximal worth would be engagement with the greatest object of worth there is, and by the principle that engagement with what is of greater worth makes a life more meaningful (Principle~P2 of EXI-11-a), it would add to the degree of meaning. It would not change what meaning is. If REL-01 is right, the additional meaning is available; if it is wrong, the meaning Definition~1 describes is all there is, and it is available to everyone. Either way, the conclusion of EXI-01 stands. Volume~II, Part~D, develops the point.
\end{reply}

\subsection{What the text does not establish}

EXI-01-a, read with the Commentary, does not establish five things that might be expected of a text with its title.

It does not say what human life as such is for. The comment on \Lm{1} has explained why the second question, what makes a life meaningful, does not depend on the answer to the first.

It does not say which projects have objective worth. It relies on the theory of the good of ETH-01 for that, and its availability claim depends only on there being, within every agent's reach, at least one.

It does not say how meaningful any particular life is, or how meaningful a life any particular agent can attain. These depend on what is done with the life and on the circumstances in which it is done, and Remark~5 sets them aside.

It does not refute the sceptic who demands that worth be vindicated from the detached standpoint. It denies her vindication premise and accepts the absurdity that her challenge reveals, and it holds that the absurdity is compatible with meaning.

It does not settle whether beings who are not agents capable of intentional action can have meaningful lives in some other sense. Its availability claim is a claim about agents.

\subsection{Conclusion}

EXI-01-a states, on one page, that life is meaningful to the extent that one engages with projects of objective worth, and that meaning in this sense is available to everyone. The Commentary has shown that the page is exact. Its definition is Wolf's, read so as to include love, receptive engagement and the work of those whose engagement is not felt as absorption. Its notion of objective worth is restricted, requiring independence only from the evaluative attitudes of those engaged, and it is entailed by the theory of the good on which the collection rests. Its theorem follows from its third assumption and a bridge principle, not from its second, whose role is to fix what the theorem means and where the text stands on the sceptical challenge. Its remark identifies the premise of that challenge that the text denies and correctly declines to claim a refutation. And its conclusion is compatible with the absurdity of life in Nagel's sense, which is real but does not touch meaning. Meaning, on this account, is the meeting of engagement and worth, and since engagement is within every agent's power and worth is present in every human situation, it is available to every agent who can engage, for as long as she can.
